% Hauptdatei. Mit XeLaTeX oder LuaLaTeX zweimal kompilieren.
\documentclass[11pt,a4paper,parskip=half]{scrartcl}
\usepackage{leipzig-handout}
% Diese Angaben fuer jedes neue Handout anpassen.
\newcommand{\HandoutTitel}{Titel des Handouts}
\newcommand{\HandoutUntertitel}{Untertitel oder Leitfrage der Sitzung}
\newcommand{\HandoutKurztitel}{Kurztitel der Sitzung}
\newcommand{\HandoutAutor}{Vorname Nachname}
\newcommand{\HandoutInstitut}{Institut oder Arbeitsgruppe}
\newcommand{\HandoutVeranstaltung}{Lehrveranstaltung oder Modul}
\newcommand{\HandoutDatum}{Datum / Semester}


% Optional: Randmuster fuer eine sparsame Druckfassung ausschalten.
% \Randmusterfalse

\begin{document}
\handouttitel

% Die folgenden Beispiele durch eigene Inhalte ersetzen.
\section{Thema und Lernziele}
Eine kurze Einleitung ordnet das Thema ein, benennt die Leitfrage
und stellt den Bezug zur Lehrveranstaltung her.

Nach der Sitzung können die Teilnehmenden beispielsweise
\begin{itemize}
  \item einen zentralen Begriff in eigenen Worten erklären,
  \item zwei Ansätze anhand geeigneter Kriterien vergleichen und
  \item das Gelernte auf eine konkrete Aufgabe anwenden.
\end{itemize}

\section{Begriffe und Kerngedanken}
\begin{merksatz}[Zentraler Begriff]
Hier steht eine kurze Definition oder eine zentrale Aussage.
Ein anschließendes Beispiel kann den Begriff veranschaulichen.
\end{merksatz}

Im Fließtext können Sie Zusammenhänge erklären und Begriffe durch
\textbf{Fettdruck} hervorheben. Zusätzliche Angaben passen in eine
Fußnote.\footnote{Beispiel einer ergänzenden Erläuterung.}

\section{Eine Gegenüberstellung}
% Die Tabelle ist direkt in den Text eingebunden und verschiebt sich nicht.
\begin{tabularx}{\linewidth}{@{}P{.22\linewidth}Y Y@{}}
  \toprule
  \textbf{Kriterium} & \textbf{Ansatz A} & \textbf{Ansatz B} \\
  \midrule
  Ziel & Beschreibung des ersten Ziels & Beschreibung des zweiten Ziels \\
  Vorgehen & Zentrale Schritte des ersten Ansatzes & Zentrale Schritte des zweiten Ansatzes \\
  Anwendung & Ein passendes Beispiel & Ein passendes Beispiel \\
  \bottomrule
\end{tabularx}
\quelle{Beispieltabelle. Inhalte und gegebenenfalls Quellenangabe ersetzen.}

\section{Kurze Reflexion}
Welche Verbindung sehen Sie zwischen dem Thema und Ihrem bisherigen Wissen?
\notizlinien[2]

% Nur fuer die Demonstration der Fortsetzungsseite.
% Ohne \newpage fliesst der eigene Inhalt automatisch weiter.
\newpage

\section{Beispiel und Anwendung}
Beschreiben Sie hier einen Fall, eine Beobachtung oder eine Aufgabe,
auf die sich die zuvor erläuterten Begriffe anwenden lassen.
Bei mathematischen Themen können Sie Formeln ergänzen, zum Beispiel
\[
  \bar{x}=\frac{1}{n}\sum_{i=1}^{n}x_i.
\]
Dabei bezeichnet $\bar{x}$ den arithmetischen Mittelwert der $n$ Werte $x_i$.

\subsection{Arbeitsauftrag}
\begin{enumerate}
  \item Beschreiben Sie den Fall in eigenen Worten.
  \item Ordnen Sie die Beobachtungen den passenden Begriffen zu.
  \item Begründen Sie Ihre Einschätzung mit Bezug auf die Inhalte der Sitzung.
\end{enumerate}

\begin{merksatz}[Ergebnissicherung]
Halten Sie Ihre wichtigste Erkenntnis in einem Satz fest. Ergänzen Sie
anschließend eine Frage, die noch offen geblieben ist.
\end{merksatz}

\section{Raum für Notizen}
\notizlinien[6]

\section{Literatur und weiterführende Hinweise}
% Platzhalter ersetzen; hier werden keine realen Quellen behauptet.
\begin{description}[font=\normalfont\bfseries,style=nextline]
  \item[Grundlagentext]
    Nachname, Vorname (Jahr): \emph{Titel des Textes}. Verlag oder Zeitschrift.
  \item[Materialien zur Sitzung]
    Titel der Materialien, Lernplattform oder Projektadresse.
\end{description}
\quelle{Platzhalter für Ihre Quellen. Für umfangreiche Literaturlisten
können Sie ein Literaturverwaltungs-Paket ergänzen.}

\end{document}
